\chapter{关键处理与交互系统源码清单}\label{app:code}

本附录摘录仓库中最核心的四段工程实现，分别对应规则清洗分类、按题聚类 Bootstrap 重算、原生 SVG 感知均匀配色映射以及静态快照 SHA-256 互锁校验。

\section{数据清洗与规则分类核心实现}

\begin{lstlisting}[language=Python, caption={scripts/cleaning/classify\_and\_audit.py 可见输出风格分类与代理预算检查}]
def classify(row, dataset, anonymous_id):
    text = row["response_text"]
    family, limit, base = prompt_family(row["prompt_id"])
    words, chars = len(text.split()), len(text)
    steps = len(STEP.findall(text))
    if not text.strip():
        style = "empty"
    elif steps:
        style = "explicit_step_markers"
    elif BULLET.search(text):
        style = "bullet_list"
    elif len(text.strip().splitlines()) > 1:
        style = "multiline_unmarked"
    else:
        style = "singleline"
    success = row.get("response_status", row.get("status")) == "ok"
    unresolved = (
        row.get("grading_status") not in {"已判分", "已按保存的逐条复核结论判分"}
        if dataset == "reasoning" else not row["extracted_answer"].strip()
    )
    state = (
        "request_failed" if not success
        else "truncated" if yes(row["is_truncated"])
        else "extraction_unresolved" if unresolved else "scored_source"
    )
    return style, state, family, limit, words, chars
\end{lstlisting}

\section{按题聚类 Bootstrap 区间复算脚本}

\begin{lstlisting}[language=Python, caption={scripts/analysis/derive\_concise\_ci.py 聚类重采样核心循环}]
for model in sorted({m for (m, _, _) in agg}):
    qs = sorted({q for (m, p, q) in agg if m == model and p in ('BeConcise', 'DefaultCoT')})
    a = [agg[(model, 'BeConcise', q)] for q in qs]
    b = [agg[(model, 'DefaultCoT', q)] for q in qs]
    point = sum(x[0] for x in a) / sum(x[1] for x in a) - sum(x[0] for x in b) / sum(x[1] for x in b)
    rng = random.Random(seed)
    diffs = []
    for _ in range(draws):
        idx = [rng.randrange(len(qs)) for _ in qs]
        diffs.append(
            sum(a[k][0] for k in idx) / sum(a[k][1] for k in idx)
            - sum(b[k][0] for k in idx) / sum(b[k][1] for k in idx)
        )
    diffs.sort()
    lo, hi = diffs[int(0.025 * draws)] * 100, diffs[int(0.975 * draws) - 1] * 100
    out[model] = [round(lo, 2), round(hi, 2)]
\end{lstlisting}

\section{原生 SVG 感知均匀配色映射}

\begin{lstlisting}[language=JavaScript, caption={app/src/zhiwei.js 中 OKLCH 感知均匀顺序色与紫—橙发散色映射}]
function oklch(L, C, hue) {
  const a = C * Math.cos(hue * Math.PI / 180), b = C * Math.sin(hue * Math.PI / 180);
  const l = (L + 0.3963377774 * a + 0.2158037573 * b) ** 3;
  const m = (L - 0.1055613458 * a - 0.0638541728 * b) ** 3;
  const s = (L - 0.0894841775 * a - 1.2914855480 * b) ** 3;
  const rgb = [
    4.0767416621 * l - 3.3077115913 * m + 0.2309699292 * s,
    -1.2684380046 * l + 2.6097574011 * m - 0.3413193965 * s,
    -0.0041960863 * l - 0.7034186147 * m + 1.7076147010 * s
  ];
  return '#' + rgb.map(v => {
    v = Math.max(0, Math.min(1, v));
    v = v <= 0.0031308 ? 12.92 * v : 1.055 * v ** (1 / 2.4) - 0.055;
    return Math.round(v * 255).toString(16).padStart(2, '0');
  }).join('');
}
\end{lstlisting}

\section{静态构建防漂移互锁校验}

\begin{lstlisting}[language=Python, caption={scripts/system/build\_web\_app.py 快照与图稿 SHA-256 一致性检查}]
code_figures = ROOT / 'reports/figures/code_generation'
code_meta = json.loads((code_figures / 'figure-metadata.json').read_text())
if code_meta['source_evidence_sha256'] != hashlib.sha256(code_path.read_bytes()).hexdigest():
    raise ValueError('Code figure evidence changed; rebuild with scripts/figures/plot_code_generation.py.')
packed_code = (APP / 'data/code-features.json.gz').read_bytes()
if hashlib.sha256(packed_code).hexdigest() != code['features_sha256']:
    raise ValueError('Code scores and evidence belong to different snapshots.')
\end{lstlisting}
